%% sec5_empirical.tex -- Section 5: Empirical Study

\section{Empirical Study: The Brazilian Lotof\'{a}cil}
\label{sec:empirical}

The \emph{Lotof\'{a}cil} (``Easy Lotto'') is a fixed-odds lottery operated
by Caixa Econ\^{o}mica Federal, Brazil, since November 2003
\citep{CEF2024}.  On each draw,
fifteen distinct numbers are selected uniformly at random from
$[25] = \{1,\ldots,25\}$; a ticket consists of exactly fifteen numbers
chosen by the player.  Prizes are awarded for matching $t \in \{11,12,13,
14,15\}$ numbers.  We study the lowest prize level $t = 11$, which occurs
most frequently.  The relevant parameters are $v=25$, $k=15$, $t=11$,
giving single-ticket prize probability $p(25,15,11) \approx 0.10589$
(equation~\eqref{eq:pone-lotofacil}), minimum pairwise overlap $m^* = 5$,
and independence upper bound $U_n = 1-(0.89411)^n$.

\subsection{Experimental design}
\label{sec:exp-design}

Three ticket-selection strategies are compared.  The \emph{random baseline}
draws $n$ tickets independently from $\binom{[25]}{15}$ with no coordination.
The \emph{greedy} strategy runs Algorithm~\ref{alg:greedy} on a pool of
$M = \max(500, 20n)$ random candidates.  The \emph{greedy+SA} strategy
starts from the greedy output and runs Algorithm~\ref{alg:sa} with
$N_{\mathrm{iter}} = \min(500 + 100n, 2000)$ iterations, initial temperature
$T_0 = 2.0$, and cooling factor $\alpha = 0.9993$.

For each strategy and value of $n$, the joint prize probability
$P_{\mathrm{win}}(\mathcal{A})$ is estimated by counting the fraction of
$N_{\mathrm{sim}} = 50{,}000$ simulated draws for which at least one ticket
achieves eleven or more matches.  A key design choice is that the same
$50{,}000$ draws are used to evaluate all three strategies in the same trial.
This paired structure removes Monte Carlo variance from the comparison:
the within-trial difference $\hat{P}_{\mathrm{SA}} - \hat{P}_{\mathrm{rand}}$
reflects ticket quality only, not random fluctuation in the draws.

For each $n$, the full pipeline is replicated for $B$ independent trials,
where $B = 50$ for $n \le 5$, $B = 40$ for $n \in \{7,10\}$, and $B = 30$
for $n \ge 15$.  The values tested are
$n \in \{1, 2, 3, 5, 7, 10, 15, 20, 25, 30, 40, 50, 75, 100\}$.
All simulations used Python~3.10 with NumPy's \texttt{default\_rng};
supplementary code is provided.

\subsection{Results}
\label{sec:results-main}

Table~\ref{tab:pwin-full} reports estimated $P_{\mathrm{win}}$ for each
strategy and value of $n$, together with the independence upper bound $U_n$
and the absolute gain of Greedy+SA over random.

\begin{table}[ht]
  \centering
  \caption{Estimated $P_{\mathrm{win}}$ (mean $\pm$ 95\% CI across $B$ trials)
           for the three strategies.  SA = Greedy+SA.  Gain is the absolute
           difference between SA and random, in percentage points.
           $U_n = 1-(0.89411)^n$.}
  \label{tab:pwin-full}
  \smallskip
  \resizebox{\linewidth}{!}{%
  \begin{tabular}{@{}rcccccc@{}}
    \toprule
    $n$ & Random & Greedy & SA & $U_n$ & Gain (pp) & $B$ \\
    \midrule
    1   & $0.1061{\pm}0.0003$ & $0.1057{\pm}0.0004$ & $0.1057{\pm}0.0004$ & 0.1059 & $-0.0$           & 50 \\
    2   & $0.1997{\pm}0.0028$ & $0.2117{\pm}0.0005$ & $0.2120{\pm}0.0005$ & 0.2006 & $+6.2$           & 50 \\
    3   & $0.2878{\pm}0.0033$ & $0.3103{\pm}0.0010$ & $0.3111{\pm}0.0010$ & 0.2852 & $+8.1$           & 50 \\
    5   & $0.4285{\pm}0.0047$ & $0.4749{\pm}0.0017$ & $0.4756{\pm}0.0018$ & 0.4286 & $+11.0^\dagger$  & 50 \\
    7   & $0.5463{\pm}0.0049$ & $0.6015{\pm}0.0027$ & $0.6025{\pm}0.0025$ & 0.5432 & $+10.3$          & 40 \\
    10  & $0.6820{\pm}0.0068$ & $0.7378{\pm}0.0025$ & $0.7388{\pm}0.0027$ & 0.6735 & $+8.3$           & 40 \\
    15  & $0.8150{\pm}0.0060$ & $0.8700{\pm}0.0034$ & $0.8717{\pm}0.0034$ & 0.8134 & $+7.0$           & 30 \\
    20  & $0.8965{\pm}0.0044$ & $0.9342{\pm}0.0018$ & $0.9360{\pm}0.0018$ & 0.8934 & $+4.4$           & 30 \\
    25  & $0.9407{\pm}0.0036$ & $0.9676{\pm}0.0012$ & $0.9677{\pm}0.0012$ & 0.9391 & $+2.9$           & 30 \\
    30  & $0.9646{\pm}0.0024$ & $0.9839{\pm}0.0006$ & $0.9841{\pm}0.0007$ & 0.9652 & $+2.0$           & 30 \\
    40  & $0.9889{\pm}0.0013$ & $0.9959{\pm}0.0002$ & $0.9961{\pm}0.0003$ & 0.9886 & $+0.7$           & 30 \\
    50  & $0.9966{\pm}0.0005$ & $0.9989{\pm}0.0001$ & $0.9990{\pm}0.0001$ & 0.9963 & $+0.2$           & 30 \\
    75  & $0.9998{\pm}0.0001$ & $1.0000{\pm}0.0000$ & $1.0000{\pm}0.0000$ & 0.9998 & $+0.02$          & 30 \\
    100 & $1.0000{\pm}0.0000$ & $1.0000{\pm}0.0000$ & $1.0000{\pm}0.0000$ & 1.0000 & ${\approx}0$     & 30 \\
    \bottomrule
  \end{tabular}}\\[2pt]
  {\footnotesize Gain: SA minus Random (percentage points).
  $\pm$: $1.96\,\mathrm{SE}$ across $B$ trials.
  $U_n=1-(0.89411)^n$.
  ${}^\dagger$ Peak absolute gain.}
\end{table}

Several features of the results deserve comment.  At $n=1$, all three
strategies produce the same value, consistent with Proposition~\ref{prop:n1}.
The gain of the optimised strategy over random peaks at $n=5$ ($+11.0$ pp)
and remains above ten percentage points at $n=7$.  For small $n$, the SA
values exceed the independence upper bound $U_n$: at $n=5$,
$P_{\mathrm{win}}^{\mathrm{SA}} = 0.476 > U_5 = 0.429$.  This is not a
paradox; it reflects the negative-correlation regime described in
Proposition~\ref{prop:mutual-excl}, where near-minimum overlap makes the
winning events approximately mutually exclusive, yielding $P_{\mathrm{win}}
\approx 2p > U_2$.  The SA refinement adds at most 0.1 percentage points
over the pure greedy solution for most values of $n$, indicating that the
greedy construction already produces near-minimum-overlap collections for
these parameters.

\subsection{Saturation}
\label{sec:saturation}

To compare the optimised strategy against its theoretical ceiling we define
coverage efficiency $\eta(n) = P_{\mathrm{win}}^{\mathrm{SA}}(n) / U_n$.
Values above one arise for the reason just noted.  As $n$ grows, $\eta(n)$
converges to one as both the optimised and random strategies approach
$P_{\mathrm{win}} = 1$.  The empirical saturation threshold, defined as the
smallest $n$ for which $P_{\mathrm{win}} > 0.999$, is approximately $n = 75$
for both strategies.  Beyond $n = 50$, the absolute gain over random falls
below $0.2$ percentage points, and the practical value of further
optimisation becomes negligible.

Table~\ref{tab:efficiency} reports $\eta(n)$ at selected values.

\begin{table}[ht]
  \centering
  \caption{Coverage efficiency $\eta(n) = P_{\mathrm{win}}^{\mathrm{SA}} / U_n$
           for the Greedy+SA strategy.  Values above one reflect
           the negative-correlation bonus of near-minimum overlap.}
  \label{tab:efficiency}
  \smallskip
  \begin{tabular}{@{}rcccc@{}}
    \toprule
    $n$ & $P_{\mathrm{win}}^{\mathrm{SA}}$ & $U_n$ & $\eta(n)$ & $U_n - P_{\mathrm{win}}^{\mathrm{SA}}$ \\
    \midrule
    2  & 0.2120 & 0.2006 & 1.057 & $-0.011$ \\
    5  & 0.4756 & 0.4286 & 1.110 & $-0.047$ \\
    10 & 0.7388 & 0.6735 & 1.097 & $-0.065$ \\
    20 & 0.9360 & 0.8934 & 1.048 & $-0.043$ \\
    30 & 0.9841 & 0.9652 & 1.020 & $-0.019$ \\
    50 & 0.9990 & 0.9963 & 1.003 & $-0.003$ \\
    75 & 1.0000 & 0.9998 & 1.000 & $\approx 0$ \\
    \bottomrule
  \end{tabular}
\end{table}
